% Generated draft: requires mathematical and source review.
\documentclass[11pt]{article}
\usepackage[margin=25mm]{geometry}
\usepackage{fontspec}
\setmainfont{Latin Modern Roman}
\usepackage{amsmath,amssymb,mathtools,unicode-math}
\setmathfont{Latin Modern Math}
\usepackage{enumitem,needspace,xurl,hyperref}
\setlength{\parindent}{0pt}
\setlength{\parskip}{5pt}
\setlength{\emergencystretch}{3em}
\newcommand{\N}{\mathbb{N}}
\newcommand{\Z}{\mathbb{Z}}
\newcommand{\Q}{\mathbb{Q}}
\newcommand{\R}{\mathbb{R}}
\newcommand{\C}{\mathbb{C}}
\newcommand{\F}{\mathbb{F}}
\newcommand{\A}{\mathbb{A}}
\newcommand{\PP}{\mathbb{P}}
\newcommand{\E}{\mathbb{E}}
\newcommand{\eps}{\varepsilon}
\newcommand{\dd}{\,\mathrm{d}}
\begin{document}
\section*{UnsolvedMath ID: 30000225}
\textbf{Source problem code:} catalogue entry
\subsection{Definitions and mathematical statement}
Let \((R,\mathfrak m)\) be a commutative Noetherian regular local ring of equal characteristic zero, meaning \(\operatorname{char}R=\operatorname{char}(R/\mathfrak m)=0\); regularity means \(\dim_{R/\mathfrak m}\mathfrak m/\mathfrak m^2=\dim R\). Let \(P\subset R\) be a prime ideal. For \(s\geq1\), define its symbolic power by
\[
 P^{(s)}=R\cap P^sR_P
 =\{f\in R:\text{there exists }u\in R\setminus P\text{ with }uf\in P^s\}.
\]
The contraction notation uses the natural localization map. The Eisenbud--Mazur conjecture, labelled Conjecture 8 in the supplied excerpt, asks whether
\[
 \forall (R,\mathfrak m)\text{ as above}\ \forall P\in\operatorname{Spec}(R),
 \qquad P^{(2)}\subseteq\mathfrak mP.
\]
Here \(\mathfrak mP\) is the ordinary ideal product. No homogeneity, positive-height, completion or finite-type hypothesis is added to the conjecture.

The surrounding original excerpt contains separate remarks and positive-characteristic results, retained here as source context rather than hypotheses or substitutes for this conjecture. Remark 7 calls its Theorem 6 an improvement of Hochster--Huneke [8, Theorem 3.7] on growth of symbolic powers in singular affine algebras, and says that the exponent \(n-1\) on a Jacobian ideal is sharp and cannot generally be lowered to \(n-2\), citing [8, Example 3.8]. Theorem 6, its parameters and that Jacobian ideal are not reproduced. The excerpt also reports counterexamples to the symbolic-square conjecture in positive and mixed characteristic, referring to [3] and [9]. These are not counterexamples under the stated equal-characteristic-zero hypotheses.

Its Theorem 9 is printed for a regular local ring of positive characteristic and a prime of height \(h>0\), with the two assertions
\[
 P^{(h+1)}\subseteq\mathfrak mP,\qquad
 R/P\text{ is }F\text{-pure}\ \Longrightarrow\ P^{(2h-1)}\subseteq P^2.
\]
In characteristic \(p>0\), \(F\)-purity ordinarily means purity of the Frobenius map \(a\mapsto a^p\), that is, injectivity after tensoring with every module. The second displayed assertion has a source-level problem at \(h=1\): for \(R=k[x]_{(x)}\) over a perfect field of characteristic \(p\), \(P=(x)\) has \(F\)-pure field quotient, but it would require \(P=P^{(1)}\subseteq P^2\). Thus the printed context needs an additional qualification or correction that cannot be recovered here. It is not asserted as a valid unrestricted theorem, and its exponent is not silently changed.

The literature triage checked 2026-08-21 reports the equal-characteristic-zero conjecture established for homogeneous and other special primes, with only special-case verification in later work and the general case still conjectural. This is dated supplied metadata. The verification field says the conjecture was isolated from theorem and remark spillover; retaining that spillover as qualified context preserves its distinctions without incorporating the apparent error into the main question.
\subsection{Short English statement}
For every prime in an equal-characteristic-zero regular local ring, must the symbolic square lie in the maximal ideal times that prime? The characteristic and symbolic-power convention are explicit. Separate source remarks on singular rings and positive characteristic are retained, including a flagged height-one inconsistency in a neighboring quoted theorem.
\subsection{Sources}
\begin{itemize}
\item Local UnsolvedMath registry, permanent ID 30000225, source code catalogue entry, ``Eisenbud--Mazur Symbolic-Square Conjecture,'' including Conjecture 8, Remark 7, Theorem 9 and the cleaned statement.
\item \emph{Kommutative Algebra (2006)}, Oberwolfach Reports, pp.\ 1073--1126, DOI 10.4171/owr/2005/19: \url{https://doi.org/10.4171/owr/2005/19}, also labelled ``original Oberwolfach Report 19/2005.'' The title uses 2006; the DOI and proposed-year field use 2005.
\item Supplied reference ``Eisenbud--Mazur, evolutions and symbolic squares'': \url{https://arxiv.org/abs/alg-geom/9511008}.
\item Supplied reference ``More, a note on the Eisenbud--Mazur conjecture'': \url{https://arxiv.org/abs/1407.5316}.
\item The original excerpt additionally cites Hochster--Huneke [8, Theorem 3.7 and Example 3.8] and source-local labels [3] and [9]. Their full numbered bibliography is absent; no identification or URL is invented for those missing entries.
\item Literature triage checked 2026-08-21; record updated 2026-09-16. No external source has been visited or independently checked here.
\end{itemize}
\end{document}
